\begin{papercard}{solove2006}
\cf{Stages} \stagetag{S} (definition of privacy problems), \stagetag{H} (harm tree built on data-flow activities), \stagetag{C} (a taxonomy usable as top-level deductive codes). An 82-page law-review article (U.\ Pa.\ L.\ Rev.\ 154:477, 2006) that replaces a single definition of ``privacy'' by a taxonomy of sixteen \emph{harmful activities} in four groups. It is purely conceptual and legal; its value for the thesis is a structured vocabulary that follows personal data from collection to dissemination, which maps well onto app data flows. The companion typology of \emph{harm types} (physical, economic, \ldots) is Citron \& Solove (2022).

\cf{Problem and method} ``Privacy is a concept in disarray. Nobody can articulate what it means.'' \pc{solove2006}{1}{Privacy is a concept in disarray. Nobody can articulate what it means}. Abstract appeals to ``privacy'' lose against concrete countervailing interests (free speech, efficiency, security) \pc{solove2006}{2}{This Article develops a taxonomy to identify privacy problems in a comprehensive and concrete manner}. Prosser (1960) identified four tort interests (intrusion, public disclosure of private facts, false light, appropriation) but only for tort law and before modern information technology \pc{solove2006}{6}{Intrusion upon the plaintiff’s seclusion or solitude, or into his private affairs}. The taxonomy is ``an attempt to identify and understand the different kinds of socially recognized privacy violations'', using law as evidence of what society recognises \pc{solove2006}{7}{The taxonomy I develop is an attempt to identify and understand the different kinds of socially recognized privacy violations}; it is not a catalogue of existing law.

\cf{Premises} Privacy is ``the relief from a range of kinds of social friction'' \pc{solove2006}{8}{Privacy is the relief from a range of kinds of social friction}; the listed activities are not inherently harmful, ``if a person consents to most of these activities, there is no privacy violation'' (with a footnote that valid consent is itself problematic, e.g.\ boilerplate) \pc{solove2006}{8}{If a person consents to most of these activities, there is no privacy violation}. The taxonomy covers only the \emph{activity} half of the equation; which matters are private is culturally variable \pc{solove2006}{8}{The full equation for a privacy violation or problem is the existence of a certain activity that causes harms or problems affecting a private matter or activity}. Privacy is a \emph{family resemblance} concept without a common essence \pc{solove2006}{9}{Elsewhere, I have argued that privacy is best understood as a family resemblance concept}. Two kinds of harm: \emph{dignitary} harms (reputation, incivility, emotional angst; Warren \& Brandeis) \pc{solove2006}{10}{The harms Warren and Brandeis spoke of are dignitary harms}, and \emph{architectural} problems, which ``involve less the overt insult or reputational harm to a person and more the creation of the risk that a person might be harmed in the future'', akin to pollution \pc{solove2006}{11}{They involve less the overt insult or reputational harm to a person and more the creation of the risk that a person might be harmed in the future}; architectural problems either enhance the risk of harm (identity theft, fraud) or upset the balance of social or institutional power (chilling effects) \pc{solove2006}{11}{Second, a particular activity can upset the balance of social or institutional power in undesirable ways}. Privacy is also constitutive of society (Simitis, Post, Cohen, Schwartz).

\cf{Structure: a model of data flowing away from the data subject} Four groups \pc{solove2006}{12}{In the taxonomy that follows, there are four basic groups of harmful activities: (1) information collection, (2) information processing, (3) information dissemination, and (4) invasion}; the model starts with the \emph{data subject}, from whom \emph{data holders} (other people, businesses, government) collect data, process it, and disseminate it; ``the general progression ... is the data moving further away from the control of the individual'' \pc{solove2006}{12}{The general progression from information collection to processing to dissemination is the data moving further away from the control of the individual}; invasions run toward the individual and need not involve information (Figure 1, p.\ 490).

\cf{The sixteen problems (definitions)}
\begin{description}[nosep,leftmargin=1.2em,labelindent=0pt]
\item[A. Information collection] \emph{Surveillance} is ``the watching, listening to, or recording of an individual's activities'' \pc{solove2006}{14}{Surveillance is the watching, listening to, or recording of an individual’s activities}; it can ensnare data beyond any originally sought and chills behaviour. \emph{Interrogation} is the pressuring of individuals to divulge information \pc{solove2006}{24}{Interrogation is the pressuring of individuals to divulge information}.
\item[B. Information processing] (data already collected) \pc{solove2006}{28}{Information processing refers to the use, storage, and manipulation of data that has been collected}: \emph{aggregation}, combining various pieces of data about a person (creates a ``digital person''; harmful even when all data is public) \pc{solove2006}{14}{Aggregation involves the combination of various pieces of data about a person}; \emph{identification}, linking information to particular individuals; \emph{insecurity}, carelessness in protecting stored information from leaks and improper access (identity theft is its overt result; ``a problem caused by the way our information is handled and protected'') \pc{solove2006}{39}{Insecurity, in short, is a problem caused by the way our information is handled and protected}; \emph{secondary use}, use of data collected for one purpose for a different purpose without consent (the purpose-specification principle) \pc{solove2006}{43}{Secondary use” is the use of data for purposes unrelated to the purposes for which the data was initially collected without the data subject’s consent}; \emph{exclusion}, failing to let the data subject know what data others hold and to participate in its handling \pc{solove2006}{14}{Exclusion concerns the failure to allow the data subject to know about the data that others have about her and participate in its handling and use}.
\item[C. Information dissemination] \emph{breach of confidentiality} (breaking a promise to keep information confidential) \pc{solove2006}{15}{Breach of confidentiality is breaking a promise to keep a person’s information confidential}; \emph{disclosure} (revealing truthful information that affects how others judge her character); \emph{exposure} (revealing nudity, grief, bodily functions); \emph{increased accessibility} (amplifying accessibility of information) \pc{solove2006}{15}{Increased accessibility is amplifying the accessibility of information}; \emph{blackmail} (threat to disclose); \emph{appropriation} (using the identity to serve another's aims); \emph{distortion} (false or misleading information) \pc{solove2006}{15}{Distortion consists of the dissemination of false or misleading information about individuals}.
\item[D. Invasion] \emph{intrusion} (invasive acts disturbing tranquility or solitude) \pc{solove2006}{15}{Intrusion concerns invasive acts that disturb one’s tranquility or solitude}; \emph{decisional interference} (government incursion into decisions about private affairs) \pc{solove2006}{15}{Decisional interference involves the government’s incursion into the data subject’s decisions regarding her private affairs}.
\end{description}
Each problem is developed with history, case law, why it is harmful, and its relation to others (e.g., insecurity relates to aggregation and identification; exclusion relates to insecurity).

\cf{Conclusion} Privacy violations ``are a group of related harms, each of which has received at least some recognition in the law'' and there is ``no common denominator'' \pc{solove2006}{82}{Privacy violations are a group of related harms, each of which has received at least some recognition in the law}. Many problems stem from beneficial activities ``much like pollution is an outgrowth of industrial production'' \pc{solove2006}{83}{many privacy problems emerge as a result of efficacious activities, much like pollution is an outgrowth of industrial production}; courts wrongly require secrecy or private places, while aggregation and increased accessibility cause problems even for public data \pc{solove2006}{83}{Aggregation creates problems even when all of the data is already available in the public domain}.

\cf{Limits (card)} Not an empirical or systematic derivation: no sampling of cases, no coding rules, no reliability; categories overlap (the author cross-references them); US law and pre-smartphone (2006), so app ecosystems, SDKs, advertising IDs, and inference by machine learning are absent; no severity or likelihood scale; assumes consent largely dissolves violation while noting that valid consent is a problem; ``socially recognised'' is argued, not measured.

\cf{Use for this thesis} (1) A mapping scheme for \emph{where} in the data flow an issue sits, usable as top-level codes for stage H and for organising findings: collection (permissions, sensors, questionnaires = surveillance, interrogation), processing (profile building, linking identifiers, SDK data, retention, reuse for advertising or model training = aggregation, identification, secondary use, insecurity, exclusion), dissemination (sharing with analytics, advertisers, employers; broken promises in policy = breach of confidentiality, increased accessibility), invasion (notifications, pressure to log, intrusive prompts). (2) The definition of secondary use and exclusion gives the vocabulary for the purpose limitation and access-right findings in policies. (3) The activity taxonomy and the harm typology of Citron \& Solove are complementary and should be linked in the harm tree (activity $\to$ harm type), not merged. (4) Use as a naming scheme only; the thesis' codes need their own operational definitions and examples, and the taxonomy should not be presented as a validated measurement instrument.
\end{papercard}
